\documentclass[11pt,a4paper]{article}
\usepackage[margin=2.5cm]{geometry}
\usepackage{amsmath,amssymb}
\usepackage{booktabs}
\usepackage{graphicx}
\usepackage{microtype}
\usepackage[hidelinks]{hyperref}
\usepackage[numbers]{natbib}

% Replace the title, authors and affiliations with your own.
\title{\textbf{Measuring Seasonal Variation in Urban Soil Moisture}}
\author{Alex Rivera\thanks{Department of Environmental Science, Example University. \texttt{alex.rivera@example.edu}}
  \and Sam Okafor\thanks{Institute for Water Research, Sample City.}}
\date{}

\newcommand{\keywords}[1]{\par\medskip\noindent\textbf{Keywords:} #1}

\begin{document}
\maketitle

\begin{abstract}
\noindent State the problem in one sentence, then say what you did, what you found and why it matters.
A good abstract is 150 to 250 words and contains no citations or undefined abbreviations.
\keywords{soil moisture; urban hydrology; remote sensing; field survey}
\end{abstract}

\section{Introduction}
\label{sec:intro}
Open with the broad context, narrow down to the specific gap in current knowledge and finish with a clear statement of your aims.
Cite prior work like this \citep{doe2021,lee2019}.

\section{Materials and methods}
\label{sec:methods}
\subsection{Study site}
Describe where and when the data were collected.

\subsection{Analysis}
Explain the statistical model. Number any equation you refer to later, such as Equation~\eqref{eq:model}.
\begin{equation}
  \theta_t = \beta_0 + \beta_1 T_t + \beta_2 P_t + \varepsilon_t
  \label{eq:model}
\end{equation}

\section{Results}
\label{sec:results}
Report the findings without interpreting them. Table~\ref{tab:summary} summarises the measurements.

\begin{table}[ht]
  \centering
  \caption{Mean volumetric soil moisture by season (replace with your data).}
  \label{tab:summary}
  \begin{tabular}{lrrr}
    \toprule
    Season & Sites & Mean (\%) & SD (\%) \\
    \midrule
    Winter & 24 & 31.2 & 4.1 \\
    Spring & 24 & 27.8 & 3.6 \\
    Summer & 23 & 18.4 & 5.2 \\
    Autumn & 24 & 25.9 & 4.4 \\
    \bottomrule
  \end{tabular}
\end{table}

\section{Discussion}
\label{sec:discussion}
Interpret the results, compare them with earlier studies and be honest about limitations.

\section{Conclusion}
Summarise the main contribution in a short paragraph and suggest next steps.

\section*{Acknowledgements}
Thank funders, collaborators and data providers here.

% Add one \bibitem per reference. The key in braces is what you pass to \citep.
\begin{thebibliography}{9}
\bibitem[Doe(2021)]{doe2021} J.~Doe, \emph{Urban Water Balance}, 2nd ed. Example Press, 2021.
\bibitem[Lee(2019)]{lee2019} K.~Lee and P.~Singh, ``Sensor networks for soil monitoring,'' \emph{Journal of Field Methods}, vol.~12, pp.~45--61, 2019.
\end{thebibliography}

\end{document}
